% ch5_d (书页 229-237 / p244-p252)
%--C-TAIL--
%JOIN%范围为
% ---- 书页 229 / p244 ----
\begin{gather*}
-\frac{\pi}{2} < v'' < +\frac{\pi}{2}\\
-\frac{\pi}{2} < u'' < +\frac{\pi}{2}\\
-\frac{\pi}{2} < v'' + u'' < \frac{\pi}{2}.
\end{gather*}

%FIG{5.16}: {Schwarzschild 时空的共形图。}

度规的 $(v'', u'')$ 部分（也就是说，在角坐标为常数的情形下）现在与 Minkowski 空间共形相关。在新坐标中，$r=0$ 处的奇点是一些直线，它们从一个渐近区域的类时无限远一直延伸到另一个渐近区域的类时无限远。

于是，极大延拓的 Schwarzschild 解的共形图看上去就像图 5.16 那样。这幅图唯一真正微妙的地方，在于必须理解 $i^+$ 和 $i^-$（未来和过去无限远）与 $r=0$ 是彼此不同的——存在大量不会撞上奇点的类时路径。正如 Kruskal 图中那样，共形图中的光锥也张开 $45^{\circ}$；主要的不同在于，整个时空被表示在一个有限的区域之内。还要注意，共形无限远的结构与 Minkowski 空间的结构一模一样，这与 Schwarzschild 时空渐近平直的论断相一致。

\section{恒星与黑洞}

我们刚刚构造的极大延拓 Schwarzschild 解讲述了一个不同凡响的故事：其中不仅有我们苦苦寻找的黑洞，还有白洞，以及另一个通过虫洞与我们宇宙相连的渐近平直区域。然而，设想这些特征在真实世界中十分常见，还为时过早。Schwarzschild 解代表的是一种高度理想化的情形：它不仅球对称，而且在整个时空中完全不含能量-动量。Birkhoff 定理意味着，球对称时空的任何真空\emph{区域}都将由 Schwarzschild 度规的\emph{一部分}来描述，但宇宙中某处物质的存在可能会戏剧性地改变整体的图像。

% ---- 书页 230 / p245 ----
一个静态的球状物体——为明确起见，就把它叫作恒星——如果半径大于 $2GM$，其外部将是 Schwarzschild 的，但不会有任何奇点或视界，而且整体结构实际上与 Minkowski 时空非常相似。当然，真实的恒星会演化，恒星有可能最终在自身引力的拉拽下坍缩，收缩到 $r=2GM$ 以下，进而落入奇点，形成黑洞。然而，这并不需要白洞，因为这样一个时空的过去与完整 Schwarzschild 解的过去毫无相似之处。描绘恒星坍缩的共形图看上去会像图 5.17 那样。内部的阴影区域不是真空，因此不由 Schwarzschild 度规描述；特别地，不存在连接到另一个宇宙的虫洞。这个时空是渐近 Minkowski 型的，只是未来的一个区域会生成事件视界。我们看到，真实的黑洞可以与极大延拓的 Schwarzschild 解共有奇点和未来视界，而不带任何白洞、过去视界或另外的渐近区域。

我们相信，这类引力坍缩绝不是恒星演化的必然终点，但在一定条件下将会发生。广义相对论对能够抵抗引力坍缩的恒星种类给出了严格的限制；对任何给定种类的物质，足够大的质量总会导致坍缩为黑洞。此外，从天体物理观测中我们已有了极好的证据，表明黑洞存在于我们的宇宙之中。

为了理解向黑洞的引力坍缩，我们应当首先理解描述球对称恒星内部的静态位形。我们不打算深入钻研这个专题，只需对内解的基本特征获得一点感觉就够了。考虑取自式 (5.11) 的一般静态、球对称度规：

%FIG{5.17}: {由坍缩恒星形成的黑洞的共形图。阴影区域包含物质，将由某个适当的动力学内解来描述；外部区域则是 Schwarzschild 的。}

% ---- 书页 231 / p246 ----
\begin{equation*}
\mathrm{d}s^2 = -e^{2\alpha(r)}\,\mathrm{d}t^2 + e^{2\beta(r)}\,\mathrm{d}r^2 + r^2\,\mathrm{d}\Omega^2.
\tag{5.133}
\end{equation*}
我们现在寻找的是非真空解，因此转向完整的 Einstein 方程，
\begin{equation*}
G_{\mu\nu} = R_{\mu\nu} - \frac{1}{2}R g_{\mu\nu} = 8\pi G\,T_{\mu\nu}.
\tag{5.134}
\end{equation*}
Einstein 张量可由 Ricci 张量 (5.14) 与标量曲率 (5.15) 得到，
\begin{align*}
G_{tt} &= \frac{1}{r^2}e^{2(\alpha-\beta)}\left(2r\,\partial_r\beta - 1 + e^{2\beta}\right)\\
G_{rr} &= \frac{1}{r^2}\left(2r\,\partial_r\alpha + 1 - e^{2\beta}\right)\\
G_{\theta\theta} &= r^2 e^{-2\beta}\left[\partial_r^2\alpha + (\partial_r\alpha)^2 - \partial_r\alpha\,\partial_r\beta + \frac{1}{r}(\partial_r\alpha - \partial_r\beta)\right]\\
G_{\phi\phi} &= \sin^2\theta\, G_{\theta\theta}.
\tag{5.135}
\end{align*}
我们把恒星本身建模为理想流体，其能动张量为
\begin{equation*}
T_{\mu\nu} = (\rho + p)U_\mu U_\nu + p\,g_{\mu\nu}.
\tag{5.136}
\end{equation*}
能量密度 $\rho$ 与压强 $p$ 将只是 $r$ 的函数。既然我们寻求的是静态解，就可以取四速度指向类时方向。归一化到 $U^\mu U_\mu = -1$，它便成为
\begin{equation*}
U_\mu = (e^{\alpha},\, 0,\, 0,\, 0),
\tag{5.137}
\end{equation*}
于是能动张量的分量为
\begin{equation*}
T_{\mu\nu} =
\begin{pmatrix}
e^{2\alpha}\rho & & & \\
& e^{2\beta}p & & \\
& & r^2 p & \\
& & & r^2(\sin^2\theta)p
\end{pmatrix}.
\tag{5.138}
\end{equation*}
这样我们便得到 Einstein 方程的三个独立分量：$tt$ 分量，
\begin{equation*}
\frac{1}{r^2}e^{-2\beta}\left(2r\,\partial_r\beta - 1 + e^{2\beta}\right) = 8\pi G\rho,
\tag{5.139}
\end{equation*}
$rr$ 分量，
\begin{equation*}
\frac{1}{r^2}e^{-2\beta}\left(2r\,\partial_r\alpha + 1 - e^{2\beta}\right) = 8\pi Gp,
\tag{5.140}
\end{equation*}

% ---- 书页 232 / p247 ----
以及 $\theta\theta$ 分量，
\begin{equation*}
e^{-2\beta}\left[\partial_r^2\alpha + (\partial_r\alpha)^2 - \partial_r\alpha\,\partial_r\beta + \frac{1}{r}(\partial_r\alpha - \partial_r\beta)\right] = 8\pi Gp.
\tag{5.141}
\end{equation*}
$\phi\phi$ 方程与 $\theta\theta$ 方程成比例，因此不必单独考虑。

我们注意到，$tt$ 方程 (5.139) 只涉及 $\beta$ 和 $\rho$。用一个新的函数 $m(r)$ 来取代 $\beta(r)$ 会带来方便，$m(r)$ 定义为
\begin{equation*}
m(r) = \frac{1}{2G}\left(r - re^{-2\beta}\right),
\tag{5.142}
\end{equation*}
或者等价地
\begin{equation*}
e^{2\beta} = \left[1 - \frac{2Gm(r)}{r}\right]^{-1},
\tag{5.143}
\end{equation*}
于是
\begin{equation*}
\mathrm{d}s^2 = -e^{2\alpha(r)}\,\mathrm{d}t^2 + \left[1 - \frac{2Gm(r)}{r}\right]^{-1}\mathrm{d}r^2 + r^2\,\mathrm{d}\Omega^2.
\tag{5.144}
\end{equation*}
度规分量 $g_{rr}$ 显然是 Schwarzschild 情形的直接推广，但 $g_{tt}$ 却不是这样。$tt$ 方程 (5.139) 变为
\begin{equation*}
\frac{dm}{dr} = 4\pi r^2\rho,
\tag{5.145}
\end{equation*}
积分可得
\begin{equation*}
m(r) = 4\pi\int_0^r \rho(r')r'^2\,\mathrm{d}r'.
\tag{5.146}
\end{equation*}
设想我们的恒星延伸到半径 $R$ 处，在此之外是真空，由 Schwarzschild 度规描述。为了使度规在这个半径处匹配，Schwarzschild 质量 $M$ 必须为
\begin{equation*}
M = m(R) = 4\pi\int_0^R \rho(r)r^2\,\mathrm{d}r.
\tag{5.147}
\end{equation*}
看起来，$m(r)$ 不过是能量密度在恒星内部的一个积分，可以解释为半径 $r$ 以内的质量。

把 $m(r)$ 解释为积分能量密度时有一个微妙之处；在一个适当的空间积分中，体积元应当是
\begin{equation*}
\sqrt{\gamma}\,\mathrm{d}^3x = e^{\beta}r^2\sin\theta\,\mathrm{d}r\,\mathrm{d}\theta\,\mathrm{d}\phi,
\tag{5.148}
\end{equation*}
其中
\begin{equation*}
\gamma_{ij}\,\mathrm{d}x^i\mathrm{d}x^j = e^{2\beta}\,\mathrm{d}r^2 + r^2\,\mathrm{d}\theta^2 + r^2\sin^2\theta\,\mathrm{d}\phi^2
\tag{5.149}
\end{equation*}

% ---- 书页 233 / p248 ----
便是空间度规。因此，真正的积分能量密度是
\begin{align*}
\bar{M} &= 4\pi\int_0^R \rho(r)r^2 e^{\beta(r)}\,\mathrm{d}r\\
&= 4\pi\int_0^R \frac{\rho(r)r^2}{\left[1 - \dfrac{2Gm(r)}{r}\right]^{1/2}}\,\mathrm{d}r.
\tag{5.150}
\end{align*}
当然，这个差别来自结合能：恒星内流体元之间的相互引力吸引产生了一种结合能，它由下式给出
\begin{equation*}
E_{\mathrm{B}} = \bar{M} - M > 0.
\tag{5.151}
\end{equation*}
结合能是把恒星中的物质驱散到无限远所需要的能量。在广义相对论中它并不总是一个有良定义的概念，但对球状恒星来说是有意义的。

用 $m(r)$ 来表示，$rr$ 方程 (5.140) 可以写成
\begin{equation*}
\frac{d\alpha}{dr} = \frac{Gm(r) + 4\pi G r^3 p}{r\left[r - 2Gm(r)\right]}.
\tag{5.152}
\end{equation*}
不直接使用 $\theta\theta$ 方程，而是转而借助能量-动量守恒 $\nabla_\mu T^{\mu\nu} = 0$，会更加方便。对我们的度规 (5.144)，容易推出 $\nu = r$ 是唯一非平庸的分量，它给出
\begin{equation*}
(\rho + p)\frac{d\alpha}{dr} = -\frac{dp}{dr}.
\tag{5.153}
\end{equation*}
把它与 (5.152) 联立，就可以消去 $\alpha(r)$，得到
\begin{equation*}
\frac{dp}{dr} = -\frac{(\rho + p)\left[Gm(r) + 4\pi G r^3 p\right]}{r\left[r - 2Gm(r)\right]}.
\tag{5.154}
\end{equation*}
这就是\textbf{Tolman--Oppenheimer--Volkoff 方程}，或简称流体静力学平衡方程。由于 $m(r)$ 通过 (5.146) 与 $\rho(r)$ 相联系，这个方程便把 $p(r)$ 与 $\rho(r)$ 联系了起来。为了得到一个封闭的方程组，我们还需要一个关系：物态方程。一般而言，它以能量密度和比熵给出压强，$p = p(\rho, S)$。我们关心的常常是熵非常小、从而可以忽略的情形；此时物态方程便取如下形式
\begin{equation*}
p = p(\rho).
\tag{5.155}
\end{equation*}
天体物理系统常常服从多方物态方程（polytropic equation of state）$p = K\rho^{\gamma}$，其中 $K$ 和 $\gamma$ 是常数。

一个简单而半现实的恒星模型来自如下假设：流体是不可压缩的——密度为常数 $\rho_*$，一直保持到恒星表面，过了表面它便为零，

% ---- 书页 234 / p249 ----
即
\begin{equation*}
\rho(r) =
\begin{cases}
\rho_*, & r < R\\
0, & r > R.
\end{cases}
\tag{5.156}
\end{equation*}
显式地给定 $\rho(r)$ 便顶替了物态方程，因为 $p(r)$ 可以由流体静力学平衡确定。接下来直接积分 (5.146) 就得到
\begin{equation*}
m(r) =
\begin{cases}
\dfrac{4}{3}\pi r^3\rho_*, & r < R\\[8pt]
\dfrac{4}{3}\pi R^3\rho_* = M, & r > R.
\end{cases}
\tag{5.157}
\end{equation*}
对流体静力学平衡方程积分，给出
\begin{equation*}
p(r) = \rho_*\left[\frac{R\sqrt{R-2GM} - \sqrt{R^3-2GMr^2}}{\sqrt{R^3-2GMr^2} - 3R\sqrt{R-2GM}}\right].
\tag{5.158}
\end{equation*}
最后我们可以由 (5.152) 得到度规分量 $g_{tt} = -e^{2\alpha(r)}$；我们发现
\begin{equation*}
e^{\alpha(r)} = \frac{3}{2}\left(1 - \frac{2GM}{R}\right)^{1/2} - \frac{1}{2}\left(1 - \frac{2GMr^2}{R^3}\right)^{1/2}, \quad r < R.
\tag{5.159}
\end{equation*}
正如所期待的那样，压强在恒星核心附近增大。实际上，对固定半径 $R$ 的恒星，如果质量超过
\begin{equation*}
M_{\mathrm{max}} = \frac{4}{9G}R.
\tag{5.160}
\end{equation*}
中心压强 $p(0)$ 就必须大于无穷大。

因此，如果我们试图把比这更大的质量挤进半径 $R$ 之内，广义相对论便不允许任何静态解；收缩到这种尺寸的恒星必然要继续收缩，最终形成一个黑洞。我们是在密度为常数这个相当强的假设下导出这一结果的，但当这一假设被大大放宽后它依然成立；\textbf{Buchdahl 定理}断言，任何合理的静态球对称内解都满足 $M < 4R/9G$。虽然严格的证明需要更多的工作，这个结果是说得通的；设想自然界存在某个最大可维持的密度，那么原则上我们所能造出的质量最大的物体将处处具有这一密度，而这正是我们考虑过的那个具体情形。

当然，这仍然不意味着现实的天体物理客体最终总会坍缩为黑洞。一颗依靠物质压强支撑的普通行星，基本上会永远存续下去（除非出现某些概率低得离谱的量子隧穿——一颗行星隧穿成某种截然不同的东西——或者最终发生质子衰变的可能性）。但大质量恒星就是另一回事了。支撑恒星的压强来自轻核聚合成较重核所产生的热量。当核燃料耗尽时，温度下降，恒星便开始在引力的作用下收缩。

% ---- 书页 235 / p250 ----
坍缩最终可能被 Fermi 简并压强止住：电子被挤得如此靠近，以至于单凭 Pauli 不相容原理（不能有两个费米子处于同一个态），它们便能抵抗进一步的压缩。靠电子简并压强支撑的恒星遗骸称为\textbf{白矮星}（white dwarf）；典型的白矮星在大小上与地球相当。质量较小的粒子在比高质量粒子更低的数密度下就会变得简并，因此核子对白矮星中的压强没有可观的贡献。白矮星是大多数恒星的最终状态，在整个宇宙中极为常见。

然而，如果总质量足够高，恒星将达到\textbf{Chandrasekhar 极限}，在那里连电子简并压强也不足以抵抗引力的拉拽。计算给出的 Chandrasekhar 极限约为 $1.4\,M_{\odot}$，其中 $M_{\odot} = 2\times 10^{33}$ g 是太阳的质量。一旦达到这一极限，恒星就被迫坍缩到更小的半径。在这一点上，电子与质子结合生成中子和中微子（逆 $\beta$ 衰变），而中微子则径直飞走。其结果是一颗\textbf{中子星}（neutron star），典型半径约为 10 km。中子星的总光度很低，但往往快速自转并具有强磁场。这种组合催生了\textbf{脉冲星}（pulsar）：它们沿着从磁极发出的喷流加速粒子，随着中子星的自转，看起来在迅速地闪烁。脉冲星由 Bell 于 1967 年发现；在短暂地猜测它们或许代表来自某个地外文明的信号之后，人们最终接受了这个更为平淡的天体物理解释。

由于中子星中心处的条件与地球上的条件非常不同，我们对物态方程的理解并不完美。尽管如此，我们相信质量足够大的中子星自身将无法抵抗引力的拉拽，并将继续坍缩；目前对中子星最大可能质量的估计在 $3$--$4\,M_{\odot}$ 左右，即\textbf{Oppenheimer--Volkoff 极限}。由于中子流体是我们所知最致密的物质（除了某些非常思辨性的设想），人们相信这种坍缩的结局是一个黑洞。

如果真有黑洞，我们会怎样知道它的存在？直接探测的根本障碍当然在于黑暗：黑洞本身不会发出任何辐射（忽略霍金辐射——那是一个非常小的效应，我们将在第 9 章讨论）。但黑洞带有极强的引力场，所以我们有望通过观测受这些场影响的物质来间接探测它们。物质落入黑洞时会升温并发出 X 射线，我们可以用卫星天文台探测到这些射线。用这种方法已经探测到大量黑洞候选体，我们宇宙中存在真实黑洞的证据是极其有力的。\footnote{关于黑洞的天体物理证据的综述，见 A.~Celotti、J.C.~Miller 与 D.W.~Sciama (1999)，\emph{Class.~Quant.~Grav.}~\textbf{16}, A3；\texttt{http://arxiv.org/abs/astro-ph/9912186}。}绝大多数候选体属于以下两类之一。一类是质量为太阳质量量级或略高一些的黑洞；它们被认为是质量非常大的恒星演化的终点。另一类描述超大质量黑洞（supermassive black holes），
% ---- 书页 236 / p251 ----
质量介于 $10^6$ 到 $10^9$ 个太阳质量之间。它们位于星系的中心，被认为是星系形成早期驱动类星体的引擎。我们自己的银河系就包含一个天体（Sgr A$^*$），据信它是一个质量至少为 $2\times 10^6\,M_{\odot}$ 的黑洞。这些超大质量黑洞形成的确切历史还不甚清楚。其他可能性还包括产生于极早期宇宙的非常小的原初黑洞（primordial black holes），以及所谓的“中量级”黑洞——质量在约一千个太阳质量量级。

物质落入黑洞时，往往会安定下来形成一个旋转的吸积盘（accretion disk），能量和角动量都随之逐渐注入洞中。因此，我们预期在天体物理情形中看到的黑洞应当是旋转的，而且观测确实与已观测到的黑洞具有非常高的自转速率相一致。在本章中，我们通过聚焦于球对称的 Schwarzschild 解排除了黑洞自转的可能性；下一章我们将转向更一般类型的黑洞。

\section{习题}

% 习题编号照原书
\textbf{1.}\quad 一只太空猴子正快乐地沿一条圆形测地线轨道绕着一个 Schwarzschild 黑洞运转。一只邪恶的狒狒在远离黑洞的地方，企图把一颗时机经过精心计算的椰子沿径向投向黑洞，好把猴子送进黑洞内部置于死地——它知道猴子无法抗拒要接住那颗下落椰子的冲动。给定猴子的质量与初始轨道半径以及椰子的质量，请说明你打算如何着手求解这个问题（但不必真的把计算做出来）。我们这位无畏的太空猴子可能落得怎样的命运？

\textbf{2.}\quad 考虑静态、圆对称的 $(2+1)$ 维时空中的一个理想流体；等价地说，即 $(3+1)$ 维中具有完全转动对称性的一个柱状位形。

\textbf{(a)}\quad 推导 $(2+1)$ 维情形下 Tolman--Oppenheimer--Volkoff（TOV）方程的类似物。

\textbf{(b)}\quad 证明真空解可以写成
\begin{equation*}
\mathrm{d}s^2 = -\mathrm{d}t^2 + \frac{1}{1-8GM}\,\mathrm{d}r^2 + r^2\,\mathrm{d}\theta^2
\end{equation*}
这里 $M$ 是常数。

\textbf{(c)}\quad 证明同一个解还有另一种写法
\begin{equation*}
\mathrm{d}s^2 = -\mathrm{d}\tau^2 + \mathrm{d}\xi^2 + \xi^2\,\mathrm{d}\phi^2
\end{equation*}
其中 $\phi\in\left[0,\,2\pi(1-8GM)^{1/2}\right]$。

\textbf{(d)}\quad 对一颗常密度恒星求解 $(2+1)$ 维 TOV 方程。求出 $p(r)$ 并解出度规。

\textbf{(e)}\quad 对具有物态方程 $p = \kappa\rho^{3/2}$ 的恒星求解 $(2+1)$ 维 TOV 方程。求出 $p(r)$ 并解出度规。

\textbf{(f)}\quad 对 (d)、(e) 两个小题中的解，求质量 $M(R) = \int_0^{2\pi}\int_0^R \rho\,\mathrm{d}r\,\mathrm{d}\theta$ 以及固有质量 $\bar{M}(R) = \int_0^{2\pi}\int_0^R \rho\sqrt{-g}\,\mathrm{d}r\,\mathrm{d}\theta$。

% ---- 书页 237 / p252 ----
\textbf{3.}\quad 考虑一个已经落入事件视界之内（$r < 2GM$）的粒子（不一定沿测地线运动）。使用通常的 Schwarzschild 坐标 $(t, r, \theta, \phi)$。试证明径向坐标必须以下式给出的最小速率减小：
\begin{equation*}
\left|\frac{\mathrm{d}r}{\mathrm{d}\tau}\right| \geq \sqrt{\frac{2GM}{r} - 1}.
\end{equation*}
计算一个粒子沿从 $r = 2GM$ 到 $r = 0$ 的轨迹所能具有的最大寿命。对质量以太阳质量计的黑洞，把这个结果用秒表示出来。试证明这一最大固有时是通过 $E \to 0$ 的自由下落实现的。

\textbf{4.}\quad 考虑真空中的 Einstein 方程，但带有一个宇宙学常数：$G_{\mu\nu} + \Lambda g_{\mu\nu} = 0$。

\textbf{(a)}\quad 求出最一般的球对称度规，所用的 $(t, r)$ 坐标在 $\Lambda = 0$ 时退化为通常的 Schwarzschild 坐标。

\textbf{(b)}\quad 像式 (5.66) 那样，用有效势写出径向测地线的运动方程。概略画出有质量粒子的有效势。

\textbf{5.}\quad 考虑一个共动观者，静止在质量为 $M$ 的 Schwarzschild 黑洞周围恒定的空间坐标 $(r_*, \theta_*, \phi_*)$ 处。该观者把一个信标投入黑洞（径直向下，沿径向轨迹）。信标以恒定波长 $\lambda_{\mathrm{em}}$（在信标静止系中）发出辐射。

\textbf{(a)}\quad 计算信标的坐标速度 $\mathrm{d}r/\mathrm{d}t$，表示为 $r$ 的函数。

\textbf{(b)}\quad 计算信标的固有速度。也就是说，设想在固定的 $r$ 处有一个共动观者，在信标经过时建立起一个局部惯性坐标系，试计算这个共动观者所测得的速度。它在 $r = 2GM$ 处是多少？

\textbf{(c)}\quad 计算由 $r_*$ 处的观者测得的波长 $\lambda_{\mathrm{obs}}$，表示为辐射被发射处半径 $r_{\mathrm{em}}$ 的函数。

\textbf{(d)}\quad 计算信标在半径 $r_{\mathrm{em}}$ 处发出的一束辐射将于 $r_*$ 处被观测到的时刻 $t_{\mathrm{obs}}$。

\textbf{(e)}\quad 试证明在很晚的时刻，红移按指数增长：$\lambda_{\mathrm{obs}}/\lambda_{\mathrm{em}} \propto e^{t_{\mathrm{obs}}/T}$。给出时间常数 $T$ 用黑洞质量 $M$ 表示的表达式。

% 第5章正文至此结束（原书书页 237 末；书页 238 起为第 6 章）。
